% ============================================================
%  Part III — Adjectives
% ============================================================
\gpart{IV}{Adjectives}

% ------------------------------------------------------------
\gsec{Which ending pattern?}

\gintro{An adjective takes an ending when it stands directly in front of a noun.
The ending depends on the type of determiner preceding it, and on the gender,
number and case of the noun.}

\begin{gtbl}{colspec={X[0.85,l] X[1.4,l] X[1.35,l]}}
  Pattern & Comes after & Example \\
  Weak   & a der-word \eng{der, dieser, jeder, welcher, alle} & der rot\en{e} Wagen \\
  Mixed  & an ein-word \eng{ein, kein, mein, ihr, unser} & ein rot\en{er} Wagen \\
  Strong & nothing at all \eng{no determiner} & rot\en{er} Wagen \\
\end{gtbl}

\begin{gnote}
A bare number is not a determiner, so the adjective takes the strong
pattern: \textit{zwei \textbf{rote} Wagen}.
\end{gnote}


\tcap{No ending needed}

An adjective takes no ending at all in these two positions.

\begin{flattbl}{colspec={X[1,l] X[2.1,l]}}
  Predicative & Der Wagen ist \textbf{rot}. \eng{after sein / werden / bleiben} \\
  Adverbial   & Er fährt \textbf{schnell}. \\
\end{flattbl}

\tcap{Stem changes before an ending}

These four adjectives change their stem before an ending is added.

\begin{flattbl}{colspec={X[0.85,l] X[0.9,l] X[1.5,l]}}
  hoch   & $\rightarrow$ hoh-   & ein hoh\en{er} Berg \\
  dunkel & $\rightarrow$ dunkl- & ein dunkl\en{er} Raum \\
  teuer  & $\rightarrow$ teur-  & ein teur\en{es} Auto \\
  sauer  & $\rightarrow$ saur-  & saur\en{e} Milch \\
\end{flattbl}

% ------------------------------------------------------------
\gsec{Weak endings — after a der-word}

der, die, das \textbullet{} dieser \textbullet{} jener \textbullet{} jeder
\textbullet{} welcher \textbullet{} mancher \textbullet{} solcher \textbullet{} alle

\begin{gendertbl}{}
  Weak & maskulin & feminin & neutrum & Plural \\
  Nominativ & -e  & -e  & -e  & -en \\
  Akkusativ & -en & -e  & -e  & -en \\
  Dativ     & -en & -en & -en & -en \\
  Genitiv   & -en & -en & -en & -en \\
\end{gendertbl}

\tcap{Example: \textit{gut}}

\begin{gendertbl}{}
  gut & maskulin & feminin & neutrum & Plural \\
  Nominativ & der gute Mann    & die gute Frau  & das gute Kind    & die guten Leute \\
  Akkusativ & den guten Mann   & die gute Frau  & das gute Kind    & die guten Leute \\
  Dativ     & dem guten Mann   & der guten Frau & dem guten Kind   & den guten Leuten \\
  Genitiv   & des guten Mannes & der guten Frau & des guten Kindes & der guten Leute \\
\end{gendertbl}

% ------------------------------------------------------------
\gsec{Mixed endings — after an ein-word}

ein \textbullet{} kein \textbullet{} mein, dein, sein, ihr, sein, unser, euer, ihr, Ihr

\begin{gendertbl}{}
  Mixed & maskulin & feminin & neutrum & Plural \\
  Nominativ & \en{-er}  & -e  & \en{-es}  & -en \\
  Akkusativ & -en & -e  & \en{-es}  & -en \\
  Dativ     & -en & -en & -en & -en \\
  Genitiv   & -en & -en & -en & -en \\
\end{gendertbl}

\begin{gnote}
Highlighted endings are the ones that differ from the weak pattern.
\end{gnote}


\tcap{Example: \textit{gut}}

\begin{gendertbl}{}
  gut & maskulin & feminin & neutrum & Plural \\
  Nominativ & ein guter Mann     & eine gute Frau   & ein gutes Kind     & keine guten Leute \\
  Akkusativ & einen guten Mann   & eine gute Frau   & ein gutes Kind     & keine guten Leute \\
  Dativ     & einem guten Mann   & einer guten Frau & einem guten Kind   & keinen guten Leuten \\
  Genitiv   & eines guten Mannes & einer guten Frau & eines guten Kindes & keiner guten Leute \\
\end{gendertbl}

% ------------------------------------------------------------
\gsec{Strong endings — no determiner}

\textit{nothing} \textbullet{} \textit{bare number} \textbullet{} viel \textbullet{} wenig \textbullet{} etwas \textbullet{} mehr 

\begin{gendertbl}{}
  Strong & maskulin & feminin & neutrum & Plural \\
  Nominativ & \en{-er} & -e  & \en{-es} & \en{-e}  \\
  Akkusativ & -en & -e  & \en{-es} & \en{-e}  \\
  Dativ     & \en{-em} & \en{-er} & \en{-em} & -en \\
  Genitiv   & -en & \en{-er} & -en & \en{-er} \\
\end{gendertbl}

\begin{gnote}
Highlighted endings are the ones that differ from the weak pattern.
\end{gnote}


\tcap{Example: \textit{gut}}

\begin{gendertbl}{}
  gut & maskulin & feminin & neutrum & Plural \\
  Nominativ & guter Wein   & gute Milch  & gutes Bier   & gute Weine \\
  Akkusativ & guten Wein   & gute Milch  & gutes Bier   & gute Weine \\
  Dativ     & gutem Wein   & guter Milch & gutem Bier   & guten Weinen \\
  Genitiv   & guten Weines & guter Milch & guten Bieres & guter Weine \\
\end{gendertbl}

% ------------------------------------------------------------
\gsec{Adjectival nouns and participles}

\tcap{Adjectives used as nouns}

Capitalise the adjective and give it the ending it would have had as an
adjective — weak after a der-word, mixed after an ein-word.

\begin{gtbl}{colspec={X[1.05,l] X[1.15,l] X[0.95,l]}}
  After a der-word & After an ein-word & Meaning \\
  der Deutsch\en{e}  & ein Deutsch\en{er}   & German \eng{m} \\
  die Deutsch\en{e}  & eine Deutsch\en{e}   & German \eng{f} \\
  die Deutsch\en{en} & keine Deutsch\en{en} & Germans \eng{pl} \\
  der Beamt\en{e}    & ein Beamt\en{er}     & official \eng{m} \\
  der Bekannt\en{e}  & ein Bekannt\en{er}   & acquaintance \eng{m} \\
\end{gtbl}

\begin{gnote}
After \textit{etwas}, \textit{nichts}, \textit{viel}, \textit{wenig} the
adjectival noun is neuter and strong: \textit{etwas Neu\textbf{es}},
\textit{nichts Gut\textbf{es}}, \textit{viel Interessant\textbf{es}}.
\end{gnote}

\tcap{Participles as adjectives}

\begin{flattbl}{colspec={X[0.9,l] X[1,l] X[1.5,l]}}
  Present & Infinitiv + d & das \textbf{lachende} Kind \eng{laughing} \\
  Past    & Partizip II   & das \textbf{gekochte} Ei \eng{boiled} \\
\end{flattbl}

\tcap{Ordinal numbers}

\begin{flattbl}{colspec={X[0.85,l] X[2.15,l]}}
  1st--19th & number + \textbf{-t} + ending: \textit{der zwei\textbf{te}}, \textit{das fünf\textbf{te}} \\
  20th up   & number + \textbf{-st} + ending: \textit{der zwanzig\textbf{ste}} \\
  Irregular & \textit{erste}, \textit{dritte}, \textit{siebte}, \textit{achte} \\
\end{flattbl}

% ------------------------------------------------------------
\gsec{Adverbs}

\gintro{German barely distinguishes adjectives from adverbs. The same word does
both jobs, and as an adverb it simply takes no ending: \textit{ein
\textbf{schneller} Wagen} but \textit{er fährt \textbf{schnell}}.}

\tcap{Time}

\begin{dtbl}[\footnotesize]{W{0.9}W{1.1}SW{0.9}W{1.1}}
\rowcolor{Ink} \hd{Deutsch} & \hd{English} & \hd{Deutsch} & \hd{English} \\ \hdrule
jetzt     & now           & gleich    & in a moment \\
sofort    & immediately   & bald      & soon \\
heute     & today         & gestern   & yesterday \\
morgen    & tomorrow      & übermorgen & the day after tomorrow \\
früher    & earlier       & später    & later \\
damals    & back then     & neulich   & recently \\
zuerst    & first         & dann      & then \\
danach    & afterwards    & endlich   & at last \\
schon     & already       & noch      & still \\
gerade    & just now      & wieder    & again \\
\end{dtbl}

\tcap{Frequency}

\begin{dtbl}[\footnotesize]{W{0.9}W{1.1}SW{0.9}W{1.1}}
\rowcolor{Ink} \hd{Deutsch} & \hd{English} & \hd{Deutsch} & \hd{English} \\ \hdrule
immer     & always        & meistens  & mostly \\
oft       & often         & manchmal  & sometimes \\
selten    & rarely        & nie       & never \\
täglich   & daily         & wöchentlich & weekly \\
\end{dtbl}

\tcap{Place}

\begin{dtbl}[\footnotesize]{W{0.9}W{1.1}SW{0.9}W{1.1}}
\rowcolor{Ink} \hd{Deutsch} & \hd{English} & \hd{Deutsch} & \hd{English} \\ \hdrule
hier      & here          & dort      & there \\
oben      & above         & unten     & below \\
links     & left          & rechts    & right \\
drinnen   & inside        & draußen   & outside \\
überall   & everywhere    & nirgends  & nowhere \\
geradeaus & straight on   & zurück    & back \\
\end{dtbl}

\begin{gnote}
\textit{hin} points away from the speaker and \textit{her} towards:
\textit{Komm \textbf{her}!} \emph{come here}, \textit{Geh \textbf{hin}!}
\emph{go there}. They attach to other words freely — \textit{hinaus},
\textit{herein}, \textit{wohin}, \textit{woher}.
\end{gnote}

% ------------------------------------------------------------
\gsec{Adverbs of degree and attitude}

\tcap{Degree}

\begin{dtbl}[\footnotesize]{W{0.9}W{1.1}SW{0.9}W{1.1}}
\rowcolor{Ink} \hd{Deutsch} & \hd{English} & \hd{Deutsch} & \hd{English} \\ \hdrule
sehr      & very          & ganz      & quite, entirely \\
ziemlich  & fairly        & besonders & especially \\
zu        & too           & genug     & enough \\
fast      & almost        & kaum      & hardly \\
etwas     & somewhat      & völlig    & completely \\
mehr      & more          & weniger   & less \\
nur       & only          & sogar     & even \\
\end{dtbl}

\tcap{Attitude — the speaker's comment on the whole sentence}

\begin{dtbl}[\footnotesize]{W{0.95}W{1.05}SW{0.95}W{1.05}}
\rowcolor{Ink} \hd{Deutsch} & \hd{English} & \hd{Deutsch} & \hd{English} \\ \hdrule
leider        & unfortunately & hoffentlich  & hopefully \\
vielleicht    & perhaps       & wahrscheinlich & probably \\
sicher        & certainly     & natürlich    & of course \\
zum Glück     & luckily       & anscheinend  & apparently \\
eigentlich    & actually      & übrigens     & by the way \\
\end{dtbl}

\begin{gnote}
An attitude adverb can stand in first position, and then the verb still comes
second: \textit{\textbf{Leider} \textbf{kann} ich nicht kommen.}
\end{gnote}

\tcap{Adverbs built with -weise}

\begin{dtbl}[\footnotesize]{W{1.1}W{0.9}SW{1.1}W{0.9}}
\rowcolor{Ink} \hd{Deutsch} & \hd{English} & \hd{Deutsch} & \hd{English} \\ \hdrule
normalerweise    & normally    & möglicherweise & possibly \\
glücklicherweise & fortunately & beispielsweise & for example \\
\end{dtbl}

\begin{gnote}
Because an adverb never takes an ending, a bare adjective after a verb is always
adverbial: \textit{Sie singt \textbf{schön}} means \emph{she sings beautifully},
not \emph{she is beautiful}.
\end{gnote}
